\section{Order-two Fourier fibers and the final $\sha{R8_3_1795}$ obstruction}
\label{sec:m3-pointed-followup}

The pointed-induction argument admits a finite integer refinement when
several different nonintegral dimensions occur.  We apply it to the two
order-two pointed subcategories of $\sha{R8_3_1768}$ and
$\sha{R8_3_1795}$.  The former acquires a necessary nontrivial pointed
associator; the latter is excluded.  We also record the completed
subgroup-dual calculation for the latter ring as an intermediate
necessary-condition result, not as a construction.

\subsection{Integer packing of the two twist fibers}

Let $g$ be an invertible object of order two in a pseudo-unitary category,
let $D$ be its global FP dimension, and write $I(g)=\bigoplus_Zm_ZZ$.
Put $z_Z=\FPdim(Z)$ and $\epsilon=\nu_2(g)\in\{1,-1\}$.
Lemma~\ref{lem:research-pointed-fourier}, or directly the central-induction
trace formula of \cite[Theorem~4.1]{NgIndicators}, gives
\begin{equation}\label{eq:pointed-c2-masses}
 \theta_Z^2=\epsilon\quad(m_Z>0),\qquad
 \sum_{\theta_Z=\alpha}m_Zz_Z
 =\sum_{\theta_Z=-\alpha}m_Zz_Z=D/2,
 \qquad \alpha^2=\epsilon.
\end{equation}
Indeed the second-indicator identity is an equality in the triangle
inequality, and the first indicator of $g\ne\one$ vanishes.
For $\epsilon=1$ the first four powers of the two possible twists have
relative traces $(1,1,1,1)$ and $(-1,1,-1,1)$ over the real FP field.
For $\epsilon=-1$, both $i$ and $-i$ give $(0,-1,0,1)$.

\begin{lemma}[Finite Fourier-fiber packing]
\label{lem:pointed-c2-packing}
Suppose all simple FP dimensions belong to $\Z+\Z\sqrt q$, with
nonnegative coefficients in this basis, and write
$D/2=T_0+T_1\sqrt q$, with $T_0,T_1\in\NN$.
Let $G=\dim\operatorname{End}(I(g))$.
For a possible forgotten central simple $v=(v_i)$ containing $g$, put
\[
 m=v_g>0,\quad \sum_i v_i d_i=z_0+z_1\sqrt q,\quad
 (c,x,y)=(m^2,mz_0,mz_1).
\]
Let $\mathcal A$ be any finite set of such triples containing every
categorically possible one.  Define $S_a=\varnothing$ for $a<0$,
$S_0=\{(0,0)\}$, and, for $1\leq a\leq G$,
\begin{equation}\label{eq:pointed-c2-reachability}
 S_a=\left(\bigcup_{(c,x,y)\in\mathcal A}
       ((x,y)+S_{a-c})\right)
       \cap\bigl([0,T_0]\times[0,T_1]\bigr).
\end{equation}
A triple $(c,x,y)$ can occur only if, for some $a$ with $c\leq a\leq G$,
\begin{equation}\label{eq:pointed-c2-support}
 (T_0,T_1)\in S_a\cap S_{G-a},\qquad
 (T_0-x,T_1-y)\in S_{a-c}.
\end{equation}
Deleting triples that fail this condition, and repeating until the
finite set stabilizes, preserves every categorically possible triple.
\end{lemma}
\begin{proof}
Adjunction identifies $m_Z$ with the coefficient of $g$ in $F(Z)$.
Thus the two fibers contribute integers $a$ and $G-a$ to
$\sum_Zm_Z^2=G$.  Equation~\eqref{eq:pointed-c2-masses} identifies the
other two coordinate sums of each fiber with $(T_0,T_1)$.
Every summand has $c=m_Z^2\geq1$.  Nonnegative dimension coordinates
ensure that every partial sum stays within the indicated rectangle.
Induction on $a$ shows that \eqref{eq:pointed-c2-reachability} enumerates
exactly all nonnegative integer packings of the available triples
with Gram contribution $a$ and dimension coordinates in that rectangle.
Removing one occurrence of a given triple from its actual fiber proves
\eqref{eq:pointed-c2-support}.  The other fiber witnesses the complementary
Gram value.  After any deletion step, both fibers of an actual category
remain witnesses for the next step.  Induction proves the last claim.

Distinct central simples with equal forgotten objects are allowed to
contribute repeated copies of the same triple.  No bound on such
repetitions, beyond the Gram and dimension bounds, is imposed.  The
lemma is therefore a necessary condition even when restriction vectors
do not distinguish central simples.
\end{proof}

For the two applications, the source tensor gives $g=X_1$, $g^2=\one$,
$\langle I(\one),I(g)\rangle=0$, and $G=8$.  In source order the
FP dimensions are
\[
\begin{array}{c|l|l}
 & (d_0,\ldots,d_7)&D/2\\\hline
\sha{R8_3_1768}
 &(1,1,5+r,6+r,7+r,7+r,6+r,5+r),\ r=\sqrt{37}
 &222+36r\\
\sha{R8_3_1795}
 &(1,1,6+3r,7+3r,7+3r,7+3r,7+3r,6+3r),\ r=\sqrt5
 &270+120r.
\end{array}
\]
The disjointness from $I(\one)$ means that pointed restrictions occur
only among the high columns of the established central-induction domain.
The finite packing filter retains $243$ of the $255$ pointed high
columns for $\sha{R8_3_1768}$, and $1187$ of the $1476$ for
$\sha{R8_3_1795}$.  It imposes no condition on columns not containing
$g$, or on the auxiliary rotation columns.  In particular the latter
are not subjected to central-twist prime-support restrictions.

The filtered columns enter the complete inherited real S/STS and
indicator cone, with $241$ rational equality rows in either case.
The all-order relative-trace compression, rather than a new twist-order
cutoff, is unchanged.  The pre-existing domain has $31$ trace types in
each case; its compression parameters are $2664$ and $360$, respectively.
The packing lemma forgets the other Gram rows and correlations between
different columns, so surviving packings or feasible cone points do not
assert existence of a center.

\subsection{Complete certificates and their categorical consequences}

\begin{theorem}\label{thm:pointed-c2-results}
The following statements hold for general complex fusion categories.
\begin{enumerate}
\item Every categorification of $\sha{R8_3_1768}$ has the nontrivial
associator class on its pointed $C_2$ subcategory.  In particular the
object $\one+g$ cannot be a connected separable subgroup algebra.
\item The ring $\sha{R8_3_1795}$ has no categorification.
\end{enumerate}
\end{theorem}
\begin{proof}
By Theorem~\ref{thm:research-four-general-spherical}, any hypothetical
categorification has a pseudo-unitary Galois conjugate.  We first work
in that positive spherical form.  The complete admissible second-indicator
assignments in source order $1,\ldots,7$ are
\[
 (\epsilon,1,1,a,1,b,c)\quad\text{for }\sha{R8_3_1768},\qquad
 (\epsilon,a,b,c,1,1,1)\quad\text{for }\sha{R8_3_1795},
 \qquad a,b,c\in\{1,-1\}.
\]
These are the exhaustive indicator domains of the inherited cone;
the source tensor and the indicator parity constraints are rechecked
by the certificate verifier.  Lemma~\ref{lem:pointed-c2-packing} and
\eqref{eq:pointed-c2-masses} give the following exact domains and results:
\[
\begin{array}{c|r|r|r|l}
\text{ring}&\epsilon&\text{allowed columns}&\text{branches checked}&\text{result}\\\hline
\sha{R8_3_1768}&+1&10711&8&\text{all infeasible}\\
\sha{R8_3_1768}&-1&10468&1&\text{feasible relaxation}\\
\sha{R8_3_1795}&+1&59261&8&\text{all infeasible}\\
\sha{R8_3_1795}&-1&58074&8&\text{all infeasible}.
\end{array}
\]
For every infeasible branch the saved exact rational covector $y$
satisfies
\[
 y\cdot c\geq0\quad\text{for every permitted column }c,
 \qquad y\cdot b<0,
\]
where $b$ is that branch's right-hand side.  Consequently the necessary
system $\sum_c x_c c=b$, $x_c\geq0$, is impossible.  Every permitted
sparse column is checked directly, as well as by the factored scoring
oracle.  Previously valid covectors for the larger domain are rechecked
on this domain.  A floating-point linear program is used only to propose
covectors; the proof consists entirely of the displayed rational
inequalities.  Full branch coverage is checked before an entire parity
class is declared impossible.

For $\sha{R8_3_1795}$ the two rows exhaust all sixteen branches,
contradicting the pseudo-unitary Galois form.  This proves the second
assertion without using the subgroup-dual calculation below.
For $\sha{R8_3_1768}$, the canonical pointed indicator $\nu_2(g)=+1$
is impossible.  The two classes in $H^3(C_2,\C^\times)=C_2$ give
$\nu_2(g)=+1$ and $-1$, respectively, so the pointed associator must
be nontrivial.  Both cohomology classes are fixed by Galois conjugation;
this conclusion therefore holds in the original category as well.
A connected separable algebra supported on the whole subgroup would
trivialize that associator, proving the algebra obstruction.
The odd assignment $(-1,1,1,1,1,1,1)$ for $\sha{R8_3_1768}$ has an
exact nonnegative rational cone witness.  It supplies no associators
or model.  The other seven odd assignments were not searched in this
refinement; they are not needed for the stated restriction.
\end{proof}

\begin{corollary}[Free pointed action on modules]
\label{cor:pointed-r1768-free-modules}
For any hypothetical categorification of $\sha{R8_3_1768}$, the
nontrivial invertible $g$ acts freely on the simple objects of every
module category.  Every indecomposable module category consequently
has even rank at least four.
\end{corollary}
\begin{proof}
If $g$ fixed a simple module object, the isomorphisms identifying its
translates would give a scalar two-cochain whose coboundary is the
restriction of the pointed associator to $C_2$.  This would trivialize
the nontrivial class forced by Theorem~\ref{thm:pointed-c2-results}.
Thus all $g$-orbits have size two.  In rank two its action matrix would
be the transposition matrix.  Every other action matrix commutes with
it, since the source fusion ring is commutative, and so has the form
$\left(\begin{smallmatrix}a&b\\b&a\end{smallmatrix}\right)$ with
$a,b\in\NN$.  The positive module dimension vector is fixed by $g$
and hence has equal coordinates.  The FP dimension of the corresponding
source simple would therefore be $a+b$, an integer.  The source has
irrational simple dimensions, giving a contradiction.
\end{proof}

Applying Corollary~\ref{cor:pointed-r1768-free-modules} to the complete
previous induction-factor domain reduces $61$ factors on $49$ possible
algebra objects to $21$ factors on $21$ objects.  The surviving module
ranks and factor counts are $(4,4),(6,6),(8,9),(12,2)$.  Here the finite
test is exact: a restriction row and its $g$-translate must have equal
multiplicity, while a row fixed by translation must occur an even number
of times.  These conditions are necessary and sufficient for a free
involution on that row multiset; rank-two factors are then excluded by
the corollary.  Among nine previously fully determined action classes,
three survive, of ranks $8,6,6$, accounting for six of the retained
factors.  The other fifteen factors have undetermined full actions and
remain in the necessary domain.  Thus this is not a classification of
module categories.  The complete factor and action checks, including
$B^{\mathsf T}B=N_A$ and all source fusion products, are recorded in
\path{pointed_followup/r1768_modules/PROOF.md}, \path{verify.py}, and
\path{verification.json} below the research evidence directory.

The complete terminal records are \path{R8_3_1768.json} and
\path{R8_3_1795.json} in
\path{evidence/rank8_research_20261022/pointed_followup/}.
Their producer and exact replayer is \path{c2_fibres.py}; it reconstructs
the frozen cone and checks equality of the complete source metadata,
including domain hashes, before applying the new filter.  The independent
audits in \path{evidence/rank8_research_20261022/r1768_r3472/}
reconstruct the source Gram matrix, independently enumerate the packing
closure, reconstruct the full phase masks, and pair every explicit
column with every separating covector.  The full R1795 audit is
\path{R1795_COMPLETE_EXCLUSION_AUDIT.md}; its even-domain and
odd-domain replayers independently check all sixteen branches.
The odd-only snapshot
\path{R1795_pointed_associator_obstruction.json} is retained to bind
the logically earlier subgroup-dual investigation.

\subsection{The completed subgroup-dual route for $\sha{R8_3_1795}$}
\label{subsec:pointed-r1795-morita}

The odd-class certificates alone force a hypothetical
$\sha{R8_3_1795}$ category to have trivial pointed associator and hence
the actual separable algebra $A=\one+g$.  Before using the final
even-class certificates, its Morita dual can be determined completely
at the level of fusion rings.  We record this intermediate argument
because it closes a natural model route and illustrates a separate
finite-cover criterion.  It is not a claim that either category exists.

\begin{proposition}[Necessary subgroup-dual fusion ring]
\label{prop:pointed-r1795-self-ring}
Using only the odd-class obstruction above, a hypothetical complex
categorification $\mathcal C$ of $\sha{R8_3_1795}$ has subgroup algebra
$A=\one+g$, and the Morita dual $\mathcal B=A\text{-}\mathrm{bimod}(\mathcal C)$
has the same fusion ring $\sha{R8_3_1795}$.  This does not assert a tensor
equivalence $\mathcal B\simeq\mathcal C$.
\end{proposition}
\begin{proof}
The free $g$-orbits of source simples are
\[
 (0,1),\quad(2,7),\quad(3,6),\quad(4,5).
\]
Choose representatives $0,2,3,4$.  The induced right modules $X_iA$
are simple because
$\operatorname{End}_A(X_iA)=\operatorname{Hom}(X_i,X_iA)$ is
one-dimensional.  They exhaust the four simple right modules, since
any simple module occurs in an induced object.  The source ring is
commutative, so left tensoring by $g$ fixes each module isomorphism
class.  An $A$-bimodule structure is a left $C_2$-linearization.
Its scalar projective multiplier is trivializable since
$H^2(C_2,\C^\times)=0$.  Each right module therefore has exactly two
linearizations, differing by the nontrivial character.  These are the
eight simple objects of $\mathcal B$.  The usual bimodule dimension
normalization gives
\[
 \FPdim_{\mathcal B}(Y_i)
 =\frac{\FPdim_{\mathcal C}(X_iA)}{\FPdim_{\mathcal C}(A)}=d_i;
\]
see \cite[Chapter~7 and Exercise~7.16.9(i)]{EGNO}.

Let $h$ be the nontrivial character bimodule over $A$.  It is an
order-two invertible acting freely on the pairs of linearizations.
It is central in the fusion ring: the underlying object of each pair
is $X_i\oplus gX_i$, and both left and right $g$-actions exchange
these two distinct simple summands.  Conjugation by the diagonal
map $(1,-1)$ changes both action signs.  Changing only the left sign
is therefore isomorphic to changing only the right sign.  This proves
equality of the left and right actions on isomorphism classes; no
half-braiding is asserted.

The action on the four right modules satisfies
\[
 (X_jA)\otimes_A Y_i\cong X_j\otimes Y_i\cong X_jX_iA
\]
as right modules.  Thus the pair quotient $h=1$ has multiplication
matrices $Q_i$ equal to the source pair quotient.  These four matrices
are linearly independent, as their unit columns form the standard basis.
Writing its nonunit basis as $a,b,c$, its products are
\begin{align*}
 a^2&=1+4a+4b+4c,& ab&=4a+5b+4c,& ac&=4a+4b+5c,\\
 b^2&=1+5a+4b+5c,& bc&=4a+5b+5c,& c^2&=1+5a+5b+4c.
\end{align*}
The respective positive dimensions are $6+3\sqrt5,7+3\sqrt5,7+3\sqrt5$.

Selfduality of the bimodule simples must be proved.  Source duality
preserves each pair.  In the signed quotient $h=-1$, write
$a_i^*=e_i a_i$ with $e_i\in\{1,-1\}$, and let $\tau$ be the
coefficient-of-unit trace.  Cyclicity and invariance under the duality
anti-involution, without any commutativity assumption, give
\[
 \tau(a_j^2a_i)=\tau((a_j^2a_i)^*)
 =e_i e_j^2\tau(a_i a_j^2)=e_i\tau(a_j^2a_i).
\]
If $e_i=-1$, this trace, and hence the signed coefficient $T_{jj}^i$,
is zero.  But $T_{jj}^i\equiv Q_{jj}^i\pmod2$.  The three odd
coefficients $Q_{22}^1=Q_{33}^2=Q_{22}^3=5$ force all $e_i=1$.
All eight simples are selfdual, so the duality anti-involution fixes
the basis and the fusion ring is commutative.

There are now ten signed variables, one per unordered triple from
$\{1,2,3\}$, with domains
\[
 -Q_{ij}^k,-Q_{ij}^k+2,\ldots,Q_{ij}^k.
\]
The unit and coefficient-of-unit entries are fixed.  Frobenius
reciprocity makes the remaining tensor symmetric in its three indices.
Changing the chosen lift of a pair changes its basis sign; the odd
variables $T_{122},T_{223},T_{233}$ transform by the independent
signs of bases $1,3,2$.  They may all be required positive, retaining
exactly one representative per sign-gauge orbit.  The complete finite
box has $2\,531\,250$ assignments.  Associativity leaves $45$ signed
tensors.  Every one expands to an eight-dimensional tensor by
\begin{equation}\label{eq:pointed-cover-reconstruction}
 N_{(i,e),(j,f)}^{(k,u)}
 =\frac{Q_{ij}^k+(-1)^{e+f+u}T_{ij}^k}{2},
 \qquad e,f,u\in\{0,1\}.
\end{equation}
Nonnegativity, the unit, duality, and all ordered associativity
identities are checked again on each full tensor.

For the remaining exclusions, apply the Galois reduction to the whole
Morita context.  Morita invariance of global and FP dimension makes
$\mathcal B$ pseudo-unitary as well; see
\cite[Theorem~2.15 and Corollary~8.14]{ENO}.
Its dimension is $D=540+240\sqrt5$, with minimal polynomial
$x^2-1080x+3600$.  If a formal codegree $f$ has monic minimal polynomial
$p$ of degree $s$, the algebraic integrality of $D/f$ implies
\[
 s\,v_\ell(3600)\geq2v_\ell(|p(0)|)
 \quad\text{for every rational prime }\ell.
\]
This follows by taking norms in a common number field, or from
\eqref{eq:research-codegree-norm}.  The exact Casimir factorizations
exclude $41$ covers this way.

Two of the four remaining covers have a formal codegree
$60+24\sqrt5$.  Central induction would then supply a simple object of
dimension $D/f=5+2\sqrt5$ by \cite[Theorem~2.13]{OstrikPiv3}.
Every object forgotten to $\mathcal B$ has radical coefficient divisible
by three, because all eight simple dimensions do.  This is impossible.
Another cover has $T=Q$; its even-labelled simples form a rank-four
fusion subcategory with dimension $270+120\sqrt5$.  This contradicts
the established general rank-four bound $24+12\sqrt3$, specifically
Corollary~\texttt{finite:cor:sharp-bounds} of the companion manuscript
\path{manuscripts/rank4/papers/companion_sections/14_finite_certification.tex}.
The sole remaining cover is identified with $\sha{R8_3_1795}$ by an
explicit based permutation, checked on all $512$ coefficients.
\end{proof}

The complete cover enumeration and witnesses are
\path{enumerate_covers.py} and \path{covers.json} in
\path{evidence/rank8_research_20261022/pointed_followup/r1795_morita/}.
The independent \path{independent_review.py} traverses the same entire
box using direct commutation of symmetric multiplication matrices,
without importing the producer's six associativity polynomials.  It
then checks every ordered identity of the full tensor and computes the
full eight-dimensional regular Casimir, reproducing the $41+2+1$
exclusions and the unique exact source match.  Its source-bound receipt
and categorical audit are \path{independent_review.json} and
\path{INDEPENDENT_REVIEW.md}.  No associator, module associator, or new
model is supplied by this calculation.

\subsection{Dimension-block screening at the seven-ring intermediate stage}
\label{subsec:pointed-dimension-block}

The following independent screen was applied to the seven rings remaining
before Theorem~\ref{thm:pointed-c2-results}.  It does not exclude any of
them; its equality cases yield further exact restrictions on possible
centers of the six surviving rings.

\begin{proposition}[Saturated dimension-block restrictions]
\label{prop:pointed-dimension-block}
Work in a pseudo-unitary Galois form of a hypothetical categorification.
For $\sha{R8_3_928}$, $\sha{R8_3_941}$, and $\sha{R8_3_982}$, every
twist-one central simple $Z$ outside the FP dimension orbit satisfies
\[
 \frac{\FPdim(Z)}{D}\in\Q(\sqrt2).
\]
Equivalently, if $\FPdim(Z)=a+b\sqrt2+c\sqrt3+d\sqrt6$, then
\begin{equation}\label{eq:pointed-quartic-dimension-mask}
 a+b=3c,\qquad a+2b=6d.
\end{equation}
For $\sha{R8_3_1768}$ every such additional twist-one central simple
has $\FPdim(Z)/D\in\Q$.  Writing its dimension as $a+b\sqrt{37}$,
this is equivalent to
\begin{equation}\label{eq:pointed-37-dimension-mask}
 6a=37b.
\end{equation}
These restrictions exclude none of the already forced constituents
of $I(\one)$.
\end{proposition}
\begin{proof}
Let $K=\Q(d_X)=\Q(D)$, let $G=\operatorname{Gal}(K/\Q)$, and put
\[
 h_\chi=\sum_{\sigma\in G}\frac{\chi(\sigma)}{\sigma(D)}.
\]
The dimension-block theorem gives
$f_\chi|h_\chi|^2\leq1$, where $f_\chi$ is the conductor of the
character $\chi$; see \cite[Lemmas~2.1 and~4.1, Theorem~5.1]{SchopierayDimensionBlock}.
The dimensional grading order is one here because $K=\Q(D)$.
We also use the equality case of the norm argument in that theorem.
Let $w_\chi$ be the signed sum of the center's dimension-orbit simples,
and let $P$ project onto the twist-one coordinates.  The coordinates
of $P S_{\mathrm{normalized}}w_\chi$ on the dimension orbit already
contribute $|G|\,|h_\chi|^2$ to its squared norm, whereas the full
norm satisfies
\[
 \|P S_{\mathrm{normalized}}w_\chi\|^2\leq |G|/f_\chi.
\]
If
$f_\chi|h_\chi|^2=1$, the operator bound in the cited proof leaves
no norm for coordinates outside that orbit.  The center row formula
therefore gives
\begin{equation}\label{eq:pointed-dimension-fourier-zero}
 \theta_Z=1,\quad Z\text{ outside the dimension orbit}
 \quad\Longrightarrow\quad
 \sum_{\sigma\in G}\chi(\sigma)
             \sigma\!\left(\frac{\FPdim(Z)}{D}\right)=0
\end{equation}
for each saturated character; compare \cite[Section~6, equation~(30)]{SchopierayDimensionBlock}.

The exact values for all seven input rings are as follows.  The final
column lists the nontrivial characters' conductor-weighted squares;
trivial-character squares are also at most one.
\[
\begin{array}{c|l|c|l}
\text{ring indices}&D&h_{\mathrm{triv}}&f_\chi|h_\chi|^2\\\hline
928,941,982&144+96\sqrt2+80\sqrt3+56\sqrt6&1/2&0,1,1\ (f_\chi=8,12,24)\\
1154,1795&540+240\sqrt5&3/10&4/9\ (f_\chi=5)\\
1768&444+72\sqrt{37}&1/6&1\ (f_\chi=37)\\
3472&700+490\sqrt2&1/7&4/25\ (f_\chi=8).
\end{array}
\]
For the quartic field,
\[
 D^{-1}=\frac18-\frac{\sqrt3}{24}-\frac{\sqrt6}{48}.
\]
The two saturated characters are the $\sqrt3$ and $\sqrt6$ Fourier
components, so \eqref{eq:pointed-dimension-fourier-zero} says that
$\FPdim(Z)/D$ lies in $\Q(\sqrt2)$.  Multiplication by the displayed
inverse gives exactly \eqref{eq:pointed-quartic-dimension-mask}.
For $\sha{R8_3_1768}$ the nontrivial Fourier sum is $-1/\sqrt{37}$.
Its saturation forces a rational quotient, equivalent to
\eqref{eq:pointed-37-dimension-mask}.

Finally, the complete formal-codegree calculation checks the forced
$ I(\one)$ constituents outside the dimension orbit.  For $928,941$
there is one remaining two-dimensional irreducible representation,
with formal codegree $4$.  For $982$ the remaining codegrees are
$16-8\sqrt2$ and $16+8\sqrt2$, each twice.  For $1768$ they are
$6$ four times and $12$ twice.  Since the corresponding center simple
has dimension ratio $1/f_E$ by \cite[Theorem~2.13]{OstrikPiv3}, all
these ratios belong to the required subfield.  Thus no forced
constituent violates the new restrictions.
\end{proof}

The source-bound exact verifier \path{verify.py} and receipt
\path{verification.json} in
\path{evidence/rank8_research_20261022/pointed_followup/dimension_block/}
reconstruct every dimension character and regular Casimir directly
from the seven canonical tensors.  They check the character conductors,
Fourier sums, field equations, and complete forced-constituent lists.
No refined cone search with these extra equations is claimed here.
The equations concern only additional twist-one simples outside the
specified dimension orbit, and only the pseudo-unitary Galois form;
they do not assert original pseudo-unitarity or unitarity.
